\section{Antecedentes de renderizado neuronal}

Antes de discutir la **destilación de puntajes**, es necesario comprender un prerrequisito clave: el **renderizado diferenciable**.

\subsection{De imágenes multivistas a reconstrucción 3D}

La tarea fundamental del renderizado neuronal consiste en: dado un conjunto de imágenes desde múltiples perspectivas de una escena, reconstruir su representación en 3D de manera que las imágenes generadas al renderizarla desde cualquier nueva perspectiva sean lo más realistas posible.

\begin{knowledgebox}{Idea central del renderizado neuronal}
El renderizado neuronal utiliza \textbf{redes neuronales} como representación de la escena en 3D. Los parámetros de la red $\theta$ codifican implícitamente la información geométrica y de apariencia de la escena. Dado un conjunto de parámetros de perspectiva de cámara $\pi$, se puede generar una imagen correspondiente en 2D mediante la función de renderizado diferenciable $g(\theta, \pi)$.

El trabajo representativo incluye NeRF (Neural Radiance Fields) y 3D Gaussian Splatting.
\end{knowledgebox}

\subsection{Objetivo de optimización de la reconstrucción 3D tradicional}

Dados $N$ imágenes de referencia desde perspectivas conocidas $\{I_i\}_{i=1}^{N}$, el objetivo de optimización para la reconstrucción en 3D es muy directo:

$$
\theta^* = \arg\min_{\theta} \sum_{i=1}^{N} \| g(\theta, \pi_i) - I_i \|_2^2
$$

Donde:
\begin{itemize}
    \item $\theta$: parámetros de la representación en 3D (por ejemplo, los pesos de la red en NeRF)
    \item $\pi_i$: parámetros de cámara de la $i$-ésima imagen (posición y orientación)
    \item $g(\theta, \pi_i)$: imagen renderizada desde la perspectiva $\pi_i$
    \item $I_i$: imagen de referencia real desde la misma perspectiva
\end{itemize}

\begin{warningbox}{La calidad de reconstrucción depende del número de imágenes de entrada}
El renderizado neuronal tradicional requiere un gran número de imágenes de entrada (generalmente $\sim$100) para obtener una reconstrucción en 3D de alta calidad. Cuando el número de imágenes disminuye a 10, 5 e incluso 1, la calidad de reconstrucción cae drásticamente: las áreas no cubiertas por los ángulos de visión presentan artefactos evidentes (como regiones negras o desenfocadas).
\end{warningbox}

\subsection{Desafíos en aplicaciones prácticas}

El renderizado neuronal enfrenta diversos desafíos en aplicaciones prácticas:

\begin{itemize}
    \item \textbf{Velocidad de reconstrucción}: gran carga computacional, no adecuado para aplicaciones en tiempo real
    \item \textbf{Velocidad de renderizado}: se requiere un renderizado rápido para soportar navegación interactiva
    \item \textbf{Escala de la escena}: desde objetos individuales hasta escenas a nivel urbano
    \item \textbf{Objetos dinámicos}: cómo manejar objetos en movimiento dentro de la escena
    \item \textbf{Número de imágenes de entrada}: cómo lograr una reconstrucción de alta calidad con el mínimo número posible de imágenes
\end{itemize}

Especialmente en escenarios de comercio electrónico (como la exhibición de productos en 3D para compras en línea), los usuarios esperan poder generar modelos 3D rápidamente a partir de solo unas pocas fotografías. Esto exige una reconstrucción de alta calidad con un número mínimo de imágenes de entrada.

\subsection{Resumen de la sección}

El renderizado neuronal proporciona un puente desde imágenes 2D hacia representaciones en 3D, cuyo núcleo es la \textbf{función de renderizado diferenciable}. Los métodos tradicionales dependen de grandes cantidades de imágenes de entrada; sin embargo, al incorporar conocimientos previos externos (como el conocimiento de modelos preentrenados de difusión), se puede reducir significativamente el número de imágenes necesarias e incluso lograr generación 3D con cero ejemplos.

\newpage

